\documentclass[11pt, a4paper]{article}
\usepackage[T1]{fontenc}
\usepackage[utf8]{inputenc}
\usepackage[french]{babel}
\usepackage[top=2cm, bottom=2cm, left=2.2cm, right=2.2cm]{geometry}
\usepackage{graphicx}
\usepackage{xcolor}
\usepackage{listings}
\usepackage{fancyhdr}
\usepackage{hyperref}
\usepackage{array}
\usepackage{colortbl}
\usepackage{booktabs}
\usepackage{float}
\usepackage{tabularx}
\usepackage{tcolorbox}
\usepackage{enumitem}
\usepackage{microtype}
\usepackage{lmodern}
\usepackage{parskip}
\usepackage{titlesec}
\usepackage{lastpage}

\tcbuselibrary{skins, breakable}

% ─── Couleurs ────────────────────────────────────────────────────────────────
\definecolor{darkblue}{RGB}{20,60,120}
\definecolor{darkgreen}{RGB}{30,130,50}
\definecolor{darkorange}{RGB}{200,100,0}
\definecolor{darkred}{RGB}{160,30,30}
\definecolor{lightgray}{RGB}{248,248,248}
\definecolor{bordergray}{RGB}{210,210,210}
\definecolor{lightblue}{RGB}{232,242,255}
\definecolor{lightgreen}{RGB}{232,248,235}
\definecolor{lightorange}{RGB}{255,243,224}
\definecolor{lightred}{RGB}{255,235,235}

% ─── Titres ──────────────────────────────────────────────────────────────────
\titleformat{\section}{\large\bfseries\color{darkblue}}{}{0em}{}[\titlerule]
\titleformat{\subsection}{\normalsize\bfseries\color{darkblue!80}}{}{0em}{}
\titleformat{\subsubsection}{\normalsize\itshape\color{darkblue!60}}{}{0em}{}
\titlespacing{\section}{0pt}{10pt}{4pt}
\titlespacing{\subsection}{0pt}{7pt}{2pt}

% ─── Listings ────────────────────────────────────────────────────────────────
\lstset{
    basicstyle=\ttfamily\footnotesize,
    keywordstyle=\color{darkblue}\bfseries,
    commentstyle=\color{gray}\itshape,
    stringstyle=\color{darkred},
    breaklines=true,
    showstringspaces=false,
    frame=single,
    rulecolor=\color{bordergray},
    backgroundcolor=\color{lightgray},
    tabsize=2,
    xleftmargin=4pt, xrightmargin=4pt,
    aboveskip=4pt, belowskip=4pt
}

% ─── Boîtes ──────────────────────────────────────────────────────────────────
\newtcolorbox{infobox}[2][]{
    colback=#2!8, colframe=#2!60!black,
    fonttitle=\bfseries\small, title={#1},
    left=5pt, right=5pt, top=3pt, bottom=3pt,
    boxrule=0.6pt
}

% ─── En-têtes ────────────────────────────────────────────────────────────────
\pagestyle{fancy}
\fancyhf{}
\lhead{\small\textcolor{darkblue}{\textbf{Eirbone P2P}}}
\rhead{\small\textcolor{darkblue}{Rapport d'état du projet}}
\cfoot{\small Page~\thepage~/~\pageref{LastPage}}
\renewcommand{\headrulewidth}{0.4pt}

% ─── Commandes utiles ────────────────────────────────────────────────────────
\newcommand{\ok}{\textcolor{darkgreen}{\textbf{OK}}}
\newcommand{\partial}{\textcolor{darkorange}{\textbf{Partiel}}}
\newcommand{\missing}{\textcolor{darkred}{\textbf{Absent}}}
\newcommand{\coche}{\textcolor{darkgreen}{$\checkmark$}}
\newcommand{\croix}{\textcolor{darkred}{$\times$}}

\hypersetup{colorlinks=true, linkcolor=darkblue, urlcolor=darkblue}

% ─────────────────────────────────────────────────────────────────────────────
\begin{document}

% ── PAGE DE TITRE ─────────────────────────────────────────────────────────────
\begin{titlepage}
\centering
\vspace*{1.5cm}

\vspace{1.8cm}

\begin{tcolorbox}[colback=lightblue, colframe=darkblue, width=0.78\textwidth,
                  title={ }, fonttitle=\bfseries]
\centering
\title{\LARGE\bfseries\color{darkblue} Rapport Eirbone :\\[0.2cm]\vspace*{0.5cm}
\large Application d'échange de gros fichiers\\ en Pair-à-Pair (P2P)\\[0.5cm]\vspace*{2.5cm}
\normalsize\normalfont ENSEIRB-MATMECA -- 2\ieme année Informatique\\[0.2cm]\vspace*{0.5cm}
Équipe :\\ Adam El Mouedden,\\ Aymane Kechna\\ \& Pierre-Louis Shrestha-Salie}
\date{}
\maketitle
\end{tcolorbox}


\end{titlepage}

\tableofcontents
\newpage

% ── 1. INTRODUCTION ───────────────────────────────────────────────────────────
\section{Introduction}

\subsection{Présentation du projet}

Eirbone est un système de partage de fichiers de type BitTorrent simplifié.
Les fichiers sont découpés en pièces de 1\,024~octets, identifiés par leur
empreinte MD5, et échangés via un protocole textuel TCP (messages terminés
par \texttt{\textbackslash n}, listes encadrées par \texttt{[ ]}).

L'architecture repose sur deux composants principaux :
\begin{itemize}[nosep]
    \item \textbf{Tracker (C)} : serveur centralisé qui maintient la liste
          des fichiers disponibles et des pairs qui les possèdent.
    \item \textbf{Peer (Java)} : client capable d'annoncer ses fichiers au
          Tracker, d'en découvrir d'autres, et d'échanger directement des
          pièces avec ses pairs.
\end{itemize}

\subsection{Architecture générale}

\begin{center}
\begin{tcolorbox}[colback=lightgray, colframe=bordergray, width=0.88\textwidth]
\begin{lstlisting}[language={}]
+------------------------------------------+
|  TRACKER (C)  --  port TCP 9000          |
|  Thread pool 4 workers, max 100 peers    |
+------------------------------------------+
        |  announce / look / getfile        |
+---------------+             +---------------+
|  PEER A (Java)|             |  PEER B (Java)|
|  TCPServer    |             |  TCPServer    |
+---------------+             +---------------+
        |       communication directe        |
        +------------------------------------+
          interested / have / getpieces / data
\end{lstlisting}
\end{tcolorbox}
\end{center}

% ── 2. TRACKER ────────────────────────────────────────────────────────────────
\section{Tracker --- Implémentation en C}

\subsection{Structures de données}

Deux structures définies dans \texttt{tracker\_state.h} constituent la base
mémoire, partagée entre les threads et protégée par un mutex
(\texttt{storage\_lock()} / \texttt{storage\_unlock()}) :

\begin{lstlisting}[language=C]
typedef struct {
    char ip[INET_ADDRSTRLEN]; int listen_port; int connected;
    char seeded_keys[100][MAX_KEY]; int seed_count;
    char leeched_keys[100][MAX_KEY]; int leech_count;
    time_t last_seen;
} Peer;

typedef struct {
    char filename[MAX_FILENAME]; long length; int piece_size;
    char key[MAX_KEY];
    char owners[MAX_PEERS_PER_FILE][64]; int owner_count;
} FileEntry;
\end{lstlisting}

\subsection{Serveur TCP et thread pool}

Le serveur écoute sur le port configuré (défaut~: 9000).
Un \textbf{thread pool} évite la création d'un thread par connexion :
4~workers consomment une file d'attente circulaire de 50~tâches,
synchronisée par mutex et variable condition.
Les paramètres sont externalisés dans \texttt{config.ini}
(port, \texttt{max\_peers}~=~100, \texttt{max\_files}~=~500,
taille du pool, niveau de log).

\subsection{Protocole Tracker-Peer --- cinq commandes}

\begin{description}[leftmargin=0pt, labelindent=0pt]
\item[\texttt{announce listen <port> seed [...] leech [...]}]
      Crée ou met à jour le pair, enregistre chaque fichier seed dans
      \texttt{FileEntry}, ajoute le pair à \texttt{owners}. Réponse~: \texttt{ok}.

\item[\texttt{look [filename=nom]}]
      Parcourt tous les \texttt{FileEntry} et retourne la liste filtrée.
      Seul le critère \texttt{filename=} est supporté
      (\texttt{filesize>}/\texttt{filesize<} non implémentés).

\item[\texttt{getfile <hash>}]
      Localise un fichier par sa clé MD5 et retourne la liste
      \texttt{owners}. Réponse~: \texttt{peers <hash> [ip:port ...]}.

\item[\texttt{update seed [...] leech [...]}]
      Retire le pair de ses anciens fichiers seed, reconstruit ses listes.
      Réponse~: \texttt{ok}.

\item[\texttt{bye}]
      Retire le pair de tous les owners, le marque hors-ligne,
      ferme la connexion. Réponse~: \texttt{bye}.
\end{description}

% ── 3. TESTS TRACKER ──────────────────────────────────────────────────────────
\section{Tests et validation du Tracker}

\subsection{Résultats des tests}

\begin{table}[H]
\centering
\begin{tabular}{clll}
\toprule
\textbf{N°} & \textbf{Commande} & \textbf{Réponse obtenue} & \textbf{Statut} \\
\midrule
1 & \texttt{announce listen 12345 seed [...] leech []} & \texttt{ok}                              & \textcolor{darkgreen}{\textbf{PASS}} \\
2 & \texttt{look []}                                   & \texttt{list [file1 1000 100 hash1]}      & \textcolor{darkgreen}{\textbf{PASS}} \\
3 & \texttt{getfile hash1}                             & \texttt{peers hash1 [127.0.0.1:12345]}    & \textcolor{darkgreen}{\textbf{PASS}} \\
4 & \texttt{update seed [hash1] leech []}              & \texttt{ok}                               & \textcolor{darkgreen}{\textbf{PASS}} \\
5 & \texttt{bye}                                       & \texttt{bye}                              & \textcolor{darkgreen}{\textbf{PASS}} \\
\midrule
\multicolumn{3}{r}{\textbf{Taux de réussite}} & \textcolor{darkgreen}{\textbf{5/5 --- 100\%}} \\
\bottomrule
\end{tabular}
\caption{Tests manuels du Tracker (via \texttt{nc localhost 9000})}
\end{table}



\subsection{Fonctionnalités opérationnelles et manquantes}

\begin{table}[H]
\centering
\begin{tabular}{lc@{\hspace{1.5cm}}lc}
\toprule
\textbf{Fonctionnalité} & \textbf{État} & \textbf{Fonctionnalité} & \textbf{État} \\
\midrule
Announce seed / leech     & \textcolor{darkgreen}{\coche} OK & Critère \texttt{filesize>} dans look & \textcolor{darkred}{\croix} Absent \\
Look sans filtre          & \textcolor{darkgreen}{\coche} OK & Critère \texttt{filesize<} dans look & \textcolor{darkred}{\croix} Absent \\
Look avec \texttt{filename=} & \textcolor{darkgreen}{\coche} OK & Validation MD5 des clés         & \textcolor{darkred}{\croix} Absent \\
Getfile                   & \textcolor{darkgreen}{\coche} OK & Timeout automatique des pairs    & \textcolor{darkred}{\croix} Absent \\
Update                    & \textcolor{darkgreen}{\coche} OK & Persistance disque               & \textcolor{darkred}{\croix} Absent \\
Bye, Thread pool, Logging & \textcolor{darkgreen}{\coche} OK & & \\
\bottomrule
\end{tabular}
\caption{Bilan des fonctionnalités du Tracker}
\end{table}

% ── 4. PEER JAVA ──────────────────────────────────────────────────────────────
\section{Peer --- Implémentation en Java}

\subsection{Architecture}

\begin{table}[H]
\centering
\begin{tabularx}{0.88\textwidth}{llX}
\toprule
\textbf{Paquetage} & \textbf{Classe} & \textbf{Rôle} \\
\midrule
\texttt{client} & \texttt{TrackerClient}   & Communication TCP avec le Tracker \\
\texttt{server} & \texttt{TCPServer}       & Serveur local pour connexions P2P entrantes \\
\texttt{peer}   & \texttt{PeerClient}      & Connexions P2P sortantes \\
\texttt{peer}   & \texttt{DownloadManager} & Orchestration des téléchargements \\
\texttt{utils}  & \texttt{FileManager}     & Création, lecture, calcul MD5 \\
\texttt{utils}  & \texttt{FileInfos}       & Métadonnées (nom, taille, hash) \\
\bottomrule
\end{tabularx}
\caption{Architecture du Peer Java}
\end{table}

\subsection{Tests d'intégration}

\textbf{App.java} valide le flux complet Peer-Tracker :

\begin{lstlisting}[language=bash]
[FileManager] Fichier cree: ./shared/ajs;eoifjo;aseifj.dat (234 octets)
[TCPServer] Started on port 2000
[TrackerClient] Connecte a localhost:9000
Peers for ajs;eoifjo;aseifj => [127.0.0.1:2000]
[TCPServer] Stopped  --  BUILD SUCCESS
\end{lstlisting}

\textbf{QuickTest.java} valide les composants unitaires :

\begin{lstlisting}[language=bash]
1. FileManager  -> File key: b730e8b9f04e3ed3...  OK  |  File size: 234 bytes  OK
2. TrackerClient -> Connexion OK  |  Annonce OK
3. Peer Discovery -> Peers found  OK
4. TCPServer (port 6666) -> start OK  |  addFile OK  |  stop OK
=== All Core Components Working ===
\end{lstlisting}

\subsection{État d'avancement}

\begin{table}[H]
\centering
\begin{tabular}{lcc}
\toprule
\textbf{Fonctionnalité} & \textbf{Tracker} & \textbf{Peer} \\
\midrule
Announce / Look / Getfile / Update / Bye & \textcolor{darkgreen}{\coche} & \textcolor{darkgreen}{\coche} \\
\midrule
\texttt{interested} / \texttt{have}      & N/A & \textcolor{darkorange}{Partiel} \\
\texttt{getpieces} / \texttt{data}       & N/A & \textcolor{darkred}{\croix~Absent} \\
BufferMap Base64                          & N/A & \textcolor{darkorange}{Partiel} \\
Téléchargement multi-sources             & N/A & \textcolor{darkred}{\croix~Absent} \\
\bottomrule
\end{tabular}
\caption{État d'avancement global}
\end{table}

% ── 5. COMMUNICATION P2P ──────────────────────────────────────────────────────
\section{Communication Pair-à-Pair}

\subsection{Protocole attendu}

\begin{center}
\begin{tcolorbox}[colback=lightgray, colframe=bordergray, width=0.82\textwidth]
\begin{lstlisting}[language={}]
Peer B (Leech)                      Peer A (Seed)
     |--- interested <hash> ----------->|
     |<-- have <hash> <buffermap_b64> --|
     |--- getpieces <hash> [0 1 3] ---->|
     |<-- data <hash> [0 1024 ...]      |
     |    <octets binaires>             |
\end{lstlisting}
\end{tcolorbox}
\end{center}

Le buffermap encode en bits la possession des pièces, puis en Base64
pour la transmission textuelle
(ex.~: 8~pièces, pièces 0,2,3 possédées $\Rightarrow$ bits \texttt{10110000}
$\Rightarrow$ \texttt{tAA=}).
Le message \texttt{data} combine un en-tête texte et une charge binaire
brute (pas d'encodage supplémentaire).

\subsection{État et travail restant}

\begin{itemize}[nosep]
    \item \texttt{interested} / \texttt{have} : classes présentes,
          appels à valider en conditions réelles.
    \item \texttt{getpieces} : classe présente, tests incomplets.
    \item \texttt{data} : transfert binaire à implémenter.
    \item Buffermap Base64 : encodage/décodage à valider.
    \item Téléchargement multi-sources et gestion des déconnexions :
          non implémentés.
\end{itemize}

% ── 6. CONCLUSION ─────────────────────────────────────────────────────────────
\section{Conclusion et perspectives}

Le Tracker est fonctionnel. 
La communication Peer-Tracker est validée.
Il reste à implémenter correctement la logique de la couche P2P.



\textbf{Travail qui reste à faire}
\begin{enumerate}[nosep]
    \item Valider \texttt{TestP2P.java} complet, buffermap Base64 et message framing
    \item Implémenter \texttt{getpieces()} et le transfert binaire \texttt{data}
    \item Ajouter robustesse réseau : timeouts, retry, déconnexions
    \item Packaging, tests finaux et livraison
\end{enumerate}

\end{document}
